% Sources verifiees (2026-07-15), RE-VERIFIEES le 2026-08-25 sur les PDF originaux
% (texte extrait par pdftotext -layout, corpus de reference) :
% - Ex. 1 : citation jury verifiee VERBATIM dans Rapport_Jury_2025.pdf, section
%   "D. Commentaires du jury" / "Analyse globale des resultats et constats".
% - Ex. 2 (Urbanloop) : E2_Appliquee_Commune_2024.pdf, questions 37, 48 et 49
%   (Pmoy = 7,2 kW ; U0 = 72 V entre rails) -- conformes au sujet.
% - Ex. 3 : la norme IEC/EN 61000-2-2 "fixe un THD de courant inferieur a 25 %" :
%   verbatim du sujet CAPET SII 2024, partie B (juste avant la question B11).
%   ATTENTION : le PDF "E1_Sujet_2024" est un DOUBLON exact (meme md5) de
%   "E2_Disciplinaire_IE_2024" ; la reference correcte est donc l'epreuve ecrite
%   disciplinaire (option ingenierie electrique), et non l'epreuve 1.
% - Ex. 3, taux d'ondulation : DEFINITION RETENUE = valeur efficace de l'ondulation
%   residuelle rapportee a la valeur moyenne, tau = Uond_eff/<uS> = racine(F^2 - 1)
%   avec F le facteur de forme. PD2 : racine(pi^2/8 - 1) = 48,3 % ; P3 : 18,3 % ;
%   PD6 : 4,2 %. Cette definition (et elle seule) est utilisee dans tout l'exercice ;
%   l'ondulation CRETE-A-CRETE, distincte, est calculee et comparee en question 2.
% ----------------------------------------------------------------------------
% PASSE COMPLEMENTARITE FICHE/EXERCICES du 2026-09-02 (protocole partage) :
%  - L'ancien Ex. 4 (HYCITY, questions E8/E9/E10) etait un DOUBLON de la fiche
%    T08 : sa q.2 reproduisait QT1 etape 2 (fP du PD3 = 0,955) et sa q.3
%    reproduisait QT2 a l'identique (UM = 631 V, U2comp = 446 V, m = 1,12).
%    Il est REMPLACE par le nouvel Ex. 5, tire du MEME sujet mais des questions
%    E2, E4, E6 et E8, non traitees par la fiche : relation v1/vD1/uS du
%    commutateur "plus positif", periode TS = T1/3, ETABLISSEMENT de <uS> par
%    calcul integral, allure et valeur efficace de iD1, puis fP du P3 etabli
%    (et non applique). Le travail intellectuel demande est donc different :
%    construction des formes d'onde + integration, la ou la fiche appliquait
%    une formule donnee et calculait un rapport de transformation.
%  - AJOUT d'un exercice docTech (nouvel Ex. 4, Urbanloop 2024 Q38 a Q43) :
%    voir le bloc de sources detaille place juste avant cet exercice.
%  - Ordre revu pour une gradation monotone : 1-2-2-2-3-3.
% ----------------------------------------------------------------------------
% - Ex. 5 (HYCITY) : E1_Sujet_2026.pdf, partie E (choix 2), questions E2/E4/E6/E8
%   (E1_Sujet_2026.md l.593-666). Cahier des charges (230/400 V - 50 Hz ; sortie
%   continue 600 V ; fP > 0,9) et enonce du commutateur "plus positif" (diodes a
%   cathodes communes, charge fortement inductive modelisee par une source de
%   courant) : conformes au sujet. v1(t) = racine(2).Veff.cos(w1.t) avec
%   w1 = 100.pi rad/s : verbatim du sujet (l.640).
%   POINT OUVERT TRANCHE le 2026-09-02 (etait : 0,827 vs 0,675 pour le P3).
%   Verification a la source (E1_Sujet_2026.md l.654-664) : la formule
%   fP = (2/pi) x racine(q) x sin(pi/q), avec "q le nombre de PHASES", n'est donnee
%   par le sujet que pour le PD3 (question E9, q = 3 -> 0,955). La question E8, qui
%   porte sur le P3, ne fournit aucune formule : le candidat doit l'etablir lui-meme
%   a partir de P_A = <u_S>.I_S et de I_Deff (question E6).
%   La formule fP = (q/pi) sin(pi/q) avec q = nombre de PULSATIONS, employee ici
%   auparavant, est APOCRYPHE : elle redonne 0,955 pour q = 6 (ce qui la faisait
%   paraitre valide) mais donne 0,827 pour le P3, valeur FAUSSE.
%   Calcul exact du P3 (charge a courant constant, diodes non commandees donc
%   cos phi1 = 1) : I_Deff = I_S/racine(3) ; P_A = (3.racine(6)/2pi).V.I_S = 1,17 V.I_S ;
%   S = 3.V.I_Deff = racine(3).V.I_S ; fP = 3.racine(2)/(2.pi) = 0,675.
%   Verification independante par le facteur de distorsion (creneau 120 deg) :
%   I_1eff/I_eff = 0,3899/0,5774 = 0,675. Les deux voies concordent.
%   L'exercice fait desormais ETABLIR fP pour le P3 (0,675). La formule du sujet
%   (valable pour le PD3) n'est plus employee ici du tout, l'etude du PD3 ayant
%   ete retiree avec le doublon.
%   NB (hors perimetre de cette passe) : la fiche fiches/T08_Redresseurs.tex porte
%   encore, en QT1, l'ancienne valeur FAUSSE fP(P3) = 0,827 obtenue par la formule
%   apocryphe (q/pi).sin(pi/q). Incoherence fiche/exercices a traiter separement.
\expandafter\ifx\csname CAPETINCLUDEMODE\endcsname\relax
\documentclass[11pt,a4paper]{article}
% preamble.tex — packages et commandes communes a toutes les fiches CAPET SII IE
\input{glyphtounicode}
\pdfgentounicode=1
\usepackage{cmap}
\usepackage[T1]{fontenc}
\usepackage[utf8]{inputenc}
\usepackage[french]{babel}
% babel[french] active a \begin{document} des raccourcis de ponctuation (:;!?)
% qui cassent la syntaxe to[...] de circuitikz (erreur "+ or - expected") ;
% on les desactive, ce qui ne change que l'espacement fin avant ces signes.
\AtBeginDocument{\shorthandoff{:;!?}}
\usepackage[a4paper, top=1.6cm, bottom=1.6cm, left=1.4cm, right=1.4cm]{geometry}
\setlength{\headheight}{14.5pt}
\usepackage{amsmath, amssymb, amsfonts}
\usepackage{siunitx}
\sisetup{locale=FR, output-decimal-marker={,}, inter-unit-product=\ensuremath{{}\cdot{}},
         group-minimum-digits=4}
% Unites non fournies par siunitx, declarees ici une fois pour toutes (fiches + exercices).
% Les gardes \ifdefined evitent tout conflit avec une redeclaration locale restee dans un
% fichier : ne pas retirer les gardes, ne pas redeclarer ces unites dans les fichiers.
\DeclareSIUnit{\jour}{j}
\ifdefined\var\else\DeclareSIUnit{\var}{var}\fi   % puissance reactive
\ifdefined\vad\else\DeclareSIUnit{\vad}{VAD}\fi   % puissance deformante
\ifdefined\tr\else\DeclareSIUnit{\tr}{tr}\fi      % tour (tr/min)
\usepackage{array, booktabs, multirow, tabularx}
\usepackage{xcolor}
\usepackage[most, breakable]{tcolorbox}
\usepackage{multicol}
\usepackage{enumitem}
\usepackage{mathtools}
\usepackage{lmodern}
% Refonte visuelle 2026 (direction "moderne coloré") : Lato pour le texte,
% maths sans-serif assorties via newtxsf (charge apres amsmath/amssymb).
\usepackage[default]{lato}
\usepackage{newtxsf}
\usepackage{microtype}
\usepackage{fancyhdr}
\usepackage{tikz}
\usetikzlibrary{positioning,arrows.meta}
\usepackage{circuitikz}
\usepackage{lastpage}

% ── Couleurs par domaine (palette modernisée 2026) ────────────────────
\definecolor{colElec}{RGB}{29, 78, 216}        % Électrotechnique — bleu vif
\definecolor{colElecLight}{RGB}{219, 234, 254} % fond bleu clair
\definecolor{colAuto}{RGB}{21, 128, 61}        % Automatique — vert
\definecolor{colMeca}{RGB}{194, 65, 12}        % Mécanique — orange
\definecolor{colTherm}{RGB}{185, 28, 28}       % Thermique — rouge
\definecolor{colQT}{RGB}{232, 244, 233}        % (hérité, plus utilisé par questiontype)
\definecolor{colJury}{RGB}{255, 248, 225}      % fond jury (crème)
\definecolor{colJuryFrame}{RGB}{202, 138, 4}   % bordure/bandeau jury (ambre)
\definecolor{colGrayLight}{RGB}{245, 245, 247} % fond gris neutre
\definecolor{colSeparator}{RGB}{180, 180, 180} % gris séparateur
\definecolor{colExo}{RGB}{8, 145, 178}          % Banque d'exercices — cyan (distinct des fiches)
% ── Encre et gris de la refonte ───────────────────────────────────────
\definecolor{colInk}{RGB}{23, 32, 42}          % texte courant (quasi-noir)
\definecolor{colSlate}{RGB}{83, 98, 118}       % gris ardoise (méta, en-têtes)
\definecolor{colRule}{RGB}{203, 213, 225}      % filets clairs
\definecolor{colJuryBg}{RGB}{255, 248, 225}    % fond crème jury
\definecolor{colPresence}{RGB}{83, 98, 118}    % bloc 5 (présence dans les sujets)
\AtBeginDocument{\color{colInk}}

% ── Style encadré exercice + corrigé (banque d'exercices) ─────────────
\tcbset{
  exoBox/.style={
    breakable,
    enhanced,
    colback=colExo!8!white,
    colframe=colExo,
    colbacktitle=colExo,
    coltitle=white,
    fonttitle=\bfseries\small,
    left=4pt, right=4pt, top=3pt, bottom=3pt,
    before skip=5pt, after skip=5pt,
  }
}

% ── Bascule enonce / corrige (banque d'exercices) ─────────────────────
% \CAPETCORRIGE est defini en ligne de commande (pdflatex "\def\CAPETCORRIGE{1}\input{...}")
% pour compiler la variante corriges d'un meme fichier source ; par defaut (non defini),
% la compilation produit la variante enonces. \CAPETSHOW{enonce}{corrige} choisit la branche.
%
% TROIS MODES (refonte 2026-09-20, manuel de revision) :
%   - aucun drapeau   -> imprime l'enonce            (variante enonces, INCHANGE)
%   - \CAPETCORRIGE   -> imprime le corrige          (variante corriges, INCHANGE)
%   - \CAPETMANUEL    -> imprime l'enonce ET met le corrige de cote pour l'annexe
%
% Le mode manuel stocke chaque corrige dans une macro (\csgdef), sans ecriture
% fichier ni re-tokenisation : le contenu (circuitikz, accents, %) traverse intact.
% L'annexe les rejoue ensuite dans l'ordre via \CAPETANNEXE.
\usepackage{etoolbox}

\newcounter{CAPETexo}          % numero continu d'exercice dans le manuel
\setcounter{CAPETexo}{0}
\newcounter{CAPETannexe}       % compteur de relecture en annexe

% Titre courant de la notion en cours (pose par \CAPETNOTION), sert d'intertitre annexe
\newcommand{\CAPETnotioncode}{}
\newcommand{\CAPETnotiontitre}{}

% \CAPETNOTION{code}{titre} : ouvre une notion dans le manuel (memorise le libelle,
% et marque le rang du premier exercice de la notion pour l'intertitre d'annexe).
\newcommand{\CAPETNOTION}[2]{%
  \renewcommand{\CAPETnotioncode}{#1}%
  \renewcommand{\CAPETnotiontitre}{#2}%
  \csgdef{CAPETnotionstart\the\numexpr\value{CAPETexo}+1\relax}{#1 -- #2}%
}

% ── \degree robuste ───────────────────────────────────────────────────
% Certains fichiers d'exercices utilisent \degree. Ni le preambule ni siunitx
% ne le definissent, et \providecommand local ne suffit pas : sous hyperref,
% \degree existe deja comme commande factice (\ERROR), donc \providecommand
% s'abstient et l'appel echoue a la compilation du manuel (erreur invisible
% en compilation isolee). On le (re)definit donc inconditionnellement ici,
% apres le chargement des packages, pour tout le corpus.
\AtBeginDocument{\providecommand{\degree}{}\renewcommand{\degree}{\ensuremath{^\circ}}}

% Libelle decoratif : les fichiers d'exercices utilisent aussi \CAPETSHOW dans
% leur \fancyhead et leur bloc de titre, pour afficher "(enonces)" ou "(corriges)".
% Ces appels-la ne sont PAS des exercices : en mode manuel ils ne doivent ni
% consommer un numero, ni alimenter l'annexe. On les reconnait a leur premier
% argument, strictement egal au libelle ci-dessous.
\newcommand{\CAPETLIBELLEENONCES}{(\'enonc\'es)}

\ifdefined\CAPETMANUEL
  % --- Mode manuel : enonce en place, corrige differe vers l'annexe ---
  \newcommand{\CAPETSHOW}[2]{%
    \def\CAPETargun{#1}%
    \ifx\CAPETargun\CAPETLIBELLEENONCES
      #1%            appel decoratif : rendu tel quel, hors numerotation
    \else
      \CAPETSHOWEXO{#1}{#2}%
    \fi
  }
  \newcommand{\CAPETSHOWEXO}[2]{%
    \stepcounter{CAPETexo}%
    \csgdef{CAPETcorrige\theCAPETexo}{#2}%
    \edef\CAPETtmpnum{\theCAPETexo}%
    #1%
    \par\vspace{2pt}%
    {\footnotesize\itshape\color{colExo}%
     \hyperref[capet:cor:\CAPETtmpnum]{Corrige de l'exercice \CAPETtmpnum\ $\rightarrow$ p.~\pageref{capet:cor:\CAPETtmpnum}}%
     \label{capet:eno:\CAPETtmpnum}}%
  }
  % \CAPETANNEXE : rejoue tous les corriges captures, dans l'ordre, avec intertitres
  \newcommand{\CAPETANNEXE}{%
    \setcounter{CAPETannexe}{0}%
    \loop\ifnum\value{CAPETannexe}<\value{CAPETexo}%
      \stepcounter{CAPETannexe}%
      \edef\CAPETcurnum{\theCAPETannexe}%
      \ifcsname CAPETnotionstart\CAPETcurnum\endcsname
        \par\medskip
        {\normalsize\bfseries\color{colExo}\csuse{CAPETnotionstart\CAPETcurnum}}%
        \par\nobreak\smallskip
      \fi
      \begin{tcolorbox}[exoBox, breakable,
        title={Corrige de l'exercice \CAPETcurnum\
               \hfill\normalfont\footnotesize
               \hyperref[capet:eno:\CAPETcurnum]{enonce p.~\pageref{capet:eno:\CAPETcurnum}}}]%
      \label{capet:cor:\CAPETcurnum}%
      \csuse{CAPETcorrige\CAPETcurnum}%
      \end{tcolorbox}%
    \repeat
  }
\else
  % --- Modes historiques, comportement strictement inchange ---
  \newcommand{\CAPETSHOW}[2]{\ifdefined\CAPETCORRIGE #2\else #1\fi}
  \newcommand{\CAPETANNEXE}{}
\fi

% ── Difficulté 1 à 3 (étoiles pleines, amssymb) ───────────────────────
\newcommand{\difficulte}[1]{%
  \ifcase#1\relax
  \or $\bigstar$\;niveau 1/3%
  \or $\bigstar\bigstar$\;niveau 2/3%
  \or $\bigstar\bigstar\bigstar$\;niveau 3/3%
  \fi
}

% ── Style encadré extrait de document technique (banque d'exercices) ──
\tcbset{
  docTech/.style={
    breakable,
    enhanced,
    colback=colGrayLight,
    colframe=colSeparator,
    fonttitle=\bfseries\small\itshape,
    coltitle=black,
    colbacktitle=colGrayLight,
    left=4pt, right=4pt, top=3pt, bottom=3pt,
    before skip=4pt, after skip=4pt,
  }
}

% ── Style tcolorbox de section principale ────────────────────────────
% Refonte 2026 : bandeau titre coloré arrondi + fond très pâle, sans cadre.
\tcbset{
  ficheSection/.style={
    breakable, enhanced jigsaw,
    frame hidden, boxrule=0pt, arc=1.8mm,
    colback=#1!4!white,
    colbacktitle=#1,
    coltitle=white,
    fonttitle=\sffamily\bfseries,
    left=7pt, right=7pt, top=5pt, bottom=5pt,
    toptitle=3pt, bottomtitle=3pt,
    before skip=10pt, after skip=8pt,
  }
}

% ── Style encadré question type ──────────────────────────────────────
% Carte blanche, bandeau sombre, liseré très clair.
\tcbset{
  questiontype/.style={
    breakable, enhanced jigsaw,
    colback=white,
    frame hidden, boxrule=0pt, arc=1.8mm,
    colbacktitle=colElec!78!black,
    coltitle=white,
    fonttitle=\sffamily\bfseries\small,
    left=7pt, right=7pt, top=4pt, bottom=4pt,
    toptitle=2.5pt, bottomtitle=2.5pt,
    before skip=8pt, after skip=8pt,
    borderline={0.6pt}{0pt}{colElec!12!white},
  }
}

% ── Style encadré jury ───────────────────────────────────────────────
% Bandeau ambre, fond crème, sans cadre.
\tcbset{
  jury/.style={
    breakable, enhanced jigsaw,
    colback=colJuryBg,
    frame hidden, boxrule=0pt, arc=1.8mm,
    colbacktitle=colJuryFrame,
    coltitle=white,
    fonttitle=\sffamily\bfseries,
    left=7pt, right=7pt, top=4pt, bottom=4pt,
    toptitle=3pt, bottomtitle=3pt,
    before skip=10pt, after skip=8pt,
  }
}

% ── En-tête / pied de page communs (refonte 2026) ─────────────────────
% #1 code fiche, #2 titre court, #3 couleur domaine
\newcommand{\ficheHeaders}[3]{%
  \fancyhead[L]{\small\sffamily{\bfseries\textcolor{#3}{#1}}\hspace{0.7em}\textcolor{colSlate}{#2}}%
  \fancyhead[R]{\small\sffamily\textcolor{colSlate}{CAPET SII IE\, --- \,Session 2026}}%
  \fancyfoot[C]{\small\sffamily\textcolor{colSlate}{\thepage\,/\,\pageref{LastPage}}}%
  \fancyhead[C]{}\fancyfoot[L]{}\fancyfoot[R]{}%
}

% ── Bloc en-tête de fiche (BLOC 0, refonte 2026) ──────────────────────
% #1 code  #2 titre  #3 couleur domaine  #4 domaine  #5 ligne présence  #6 sous-thèmes
\newcommand{\ficheEntete}[6]{%
  \begin{tcolorbox}[enhanced, breakable,
    frame hidden, boxrule=0pt, arc=2.5mm,
    interior style={left color=#3, right color=#3!75!black},
    left=9pt, right=9pt, top=7pt, bottom=7pt,
    before skip=0pt, after skip=6pt]
    {\color{white}\sffamily\bfseries\LARGE #1\quad #2}\\[4pt]
    {\color{#3!22!white}\sffamily\small \MakeUppercase{#4}\quad\textbar\quad #5}\\[1pt]
    {\color{#3!35!white}\sffamily\footnotesize Sous-thèmes : #6}
  \end{tcolorbox}%
}

% ── Mise en page ─────────────────────────────────────────────────────
\pagestyle{fancy}
\fancyhf{}
\renewcommand{\headrulewidth}{0.5pt}
\renewcommand{\footrulewidth}{0pt}
\AtBeginDocument{%
  \renewcommand{\headrule}{{\color{colRule}\hrule width\headwidth height 0.5pt}}%
}

\setlist[itemize]{leftmargin=1.2em, itemsep=1pt, topsep=2pt, parsep=0pt}
\setlist[enumerate]{leftmargin=1.2em, itemsep=3pt, topsep=2pt, parsep=0pt}

\setlength{\parskip}{2pt}
\setlength{\columnsep}{10pt}

\def\CAPETBEGINFILE{\begin{document}}
\def\CAPETENDFILE{\end{document}}
\else
\def\CAPETBEGINFILE{}
\def\CAPETENDFILE{}
\fi

\fancyhead[L]{\textbf{\textcolor{colExo}{T08}} - Exercices - Redresseurs (monophas\'e et triphas\'e)}
\fancyhead[R]{\small\textit{Banque d'exercices \CAPETSHOW{(\'enonc\'es)}{(corrig\'es)}}}
\fancyfoot[C]{\small\thepage\,/\,\pageref{LastPage}}

\CAPETBEGINFILE

\begin{tcolorbox}[
  enhanced, breakable,
  colback=colExo,
  colframe=colExo,
  left=6pt, right=6pt, top=5pt, bottom=5pt,
  before skip=0pt, after skip=4pt
]
  {\color{white}\bfseries\large T08 \quad Redresseurs (monophas\'e et triphas\'e)}\\[2pt]
  {\color{white}\small Banque d'exercices -- 6 exercices \CAPETSHOW{(\'enonc\'es)}{(corrig\'es)}}
\end{tcolorbox}

% ============================================================
% EXERCICE 1 (niveau 1) -- PD2 a diodes, grandeurs moyennes et efficaces
% ============================================================
\begin{tcolorbox}[exoBox,
  title={Exercice 1 -- \difficulte{1} \hfill\normalfont\footnotesize exercice in\'edit}
]
\CAPETSHOW{
Un pont de Gra\"etz monophas\'e (PD2, 4 diodes suppos\'ees parfaites) redresse une tension alternative sinuso\"idale $v(t)$ de valeur efficace $V_{eff} = \SI{230}{\volt}$ et de fr\'equence $f = \SI{50}{\hertz}$. Il alimente directement (sans filtrage) une r\'esistance de charge $R = \SI{100}{\ohm}$.

\begin{center}
\begin{circuitikz}[european, scale=0.9]
\draw (2,2) node[circ]{} to[sV, l=$v(t)$] (5,2) node[circ]{};
\draw (2,2) to[D, l=$D_1$] (2,4.2);
\draw (5,2) to[D, l_=$D_3$] (5,4.2);
\draw (2,-0.2) to[D, l=$D_4$] (2,2);
\draw (5,-0.2) to[D, l_=$D_2$] (5,2);
\draw (2,4.2) -- (7.2,4.2);
\draw (2,-0.2) -- (7.2,-0.2);
\node[circ] at (5,4.2) {};
\node[circ] at (5,-0.2) {};
\draw (7.2,4.2) to[R, l=$R$] (7.2,-0.2);
\draw[-latex] (8.5,4.2) -- (8.5,-0.2) node[midway,right]{$u_S(t)$};
\end{circuitikz}
\end{center}

\begin{enumerate}[leftmargin=1.5em]
  \item Rappeler la relation liant la valeur moyenne $\langle u_S \rangle$ de la tension redress\'ee et $V_{eff}$ pour un pont PD2, puis calculer num\'eriquement $\langle u_S \rangle$.
  \item En d\'eduire le courant moyen $\langle i_S \rangle$ dans la r\'esistance de charge.
  \item Justifier que la valeur efficace $U_{S,eff}$ de la tension redress\'ee est \'egale \`a $V_{eff}$, puis calculer la valeur efficace $I_{S,eff}$ du courant dans $R$.
  \item Chaque diode ne conduit que pendant une moiti\'e de la p\'eriode de $v(t)$ (les diodes $D_1$-$D_2$ conduisent ensemble sur une alternance, $D_3$-$D_4$ sur l'autre). En d\'eduire la valeur efficace $I_{D,eff}$ du courant dans une diode, en fonction de $I_{S,eff}$, puis num\'eriquement.
\end{enumerate}

}{
\textbf{Corrig\'e.}

\begin{enumerate}[leftmargin=1.5em]
  \item Pour un pont PD2 :
  \[ \langle u_S \rangle = \frac{2\sqrt{2}}{\pi}\,V_{eff} \approx 0{,}9\,V_{eff} = 0{,}9 \times 230 \approx \SI{207}{\volt} \]
  \item Charge purement r\'esistive : $\langle i_S \rangle = \dfrac{\langle u_S \rangle}{R} = \dfrac{207}{100} \approx \SI{2,07}{\ampere}$.
  \item $u_S(t) = \sqrt{2}\,V_{eff}\,|\sin(\omega t)|$. Or $\sin^2$ et $|\sin|^2$ co\"incident : la valeur efficace du redressement double alternance est identique \`a celle du signal non redress\'e, soit $U_{S,eff} = V_{eff} = \SI{230}{\volt}$. Comme $R$ est lin\'eaire, $I_{S,eff} = U_{S,eff}/R = 230/100 = \SI{2,3}{\ampere}$.
  \item Sur une p\'eriode de $v(t)$, une diode (par exemple $D_1$-$D_2$) conduit le courant $i_S(t)$ pendant une demi-p\'eriode et est bloqu\'ee sur l'autre demi-p\'eriode :
  \[
  I_{D,eff} = \sqrt{\frac{1}{T}\int_0^{T/2} i_S^2(t)\,\mathrm dt} = \sqrt{\frac12}\cdot I_{S,eff} = \frac{I_{S,eff}}{\sqrt2}
  \]
  car la moyenne de $i_S^2(t)$ sur la demi-p\'eriode conductrice est \'egale \`a la moyenne de $i_S^2(t)$ sur la p\'eriode compl\`ete (m\^eme forme p\'eriodique). Num\'eriquement : $I_{D,eff} = 2{,}3/\sqrt{2} \approx \SI{1,63}{\ampere}$.
\end{enumerate}

\smallskip
\textbf{Conclusion.} Le calibre en courant des diodes se choisit sur $I_{D,eff} \approx \SI{1,63}{\ampere}$ (et non sur $I_{S,eff} = \SI{2,3}{\ampere}$) : un pont dont les diodes sont donn\'ees pour \SI{2}{\ampere} efficaces convient donc \`a cette charge, alors qu'un dimensionnement fait \`a tort sur le courant de charge conduirait \`a sur-dimensionner le composant d'un facteur $\sqrt2$.

\smallskip
\textit{Point de vigilance jury (rapport de jury CAPET 2025, \guillemotleft\ Commentaires du jury \guillemotright) : \guillemotleft\ De nombreux candidats ne maitrisent pas les lois basiques de l'\'electricit\'e (la loi des n\oe uds, la loi des mailles, les calculs de puissance ou encore la loi d'Ohm) \guillemotright\ -- une ma\^itrise s\^ure des relations valeur moyenne / valeur efficace est indispensable avant d'aborder les montages redresseurs plus complexes.}
}
\end{tcolorbox}


% ============================================================
% EXERCICE 2 (niveau 2) -- Urbanloop, reversibilite du pont a diodes
% ============================================================
\begin{tcolorbox}[exoBox,
  title={Exercice 2 -- \difficulte{2} \hfill\normalfont\footnotesize 2024, \'Epreuve \'ecrite disciplinaire appliqu\'ee, questions 37, 48 et 49}
]
\CAPETSHOW{
Le syst\`eme de transport Urbanloop alimente ses capsules \`a travers des rails aliment\'es par des bo\^itiers d'alimentation r\'epartis le long de la voie. Chaque bo\^itier est raccord\'e au r\'eseau EDF monophas\'e et impose en sortie une tension continue $U_0 = \SI{72}{\volt}$ (TBTS) entre les deux rails. La structure interne d'un bo\^itier est : pont redresseur PD2, filtre, hacheur s\'erie.

\begin{center}
\begin{tikzpicture}[scale=1, every node/.style={font=\small}]
\draw (0,0) rectangle (2.3,1) node[pos=.5] {R\'eseau EDF};
\node at (1.15,-0.35) {$\SI{230}{\volt}\sim$};
\draw[-latex] (2.3,0.5) -- (3.1,0.5);
\draw (3.1,0) rectangle (5.2,1) node[pos=.5] {Pont PD2};
\draw[-latex] (5.2,0.5) -- (6,0.5);
\draw (6,0) rectangle (7.9,1) node[pos=.5] {Filtre};
\draw[-latex] (7.9,0.5) -- (8.7,0.5);
\draw (8.7,0) rectangle (11.2,1) node[pos=.5] {Hacheur s\'erie};
\draw[-latex] (11.2,0.5) -- (12,0.5);
\draw (12,0) rectangle (13.8,1) node[pos=.5] {Rails};
\node at (12.9,-0.35) {$U_0=\SI{72}{\volt}=$};
\end{tikzpicture}
\end{center}

La puissance moyenne consomm\'ee par l'ensemble des capsules en mouvement sur le trajet est $P_{moy} = \SI{7,2}{\kilo\watt}$ (toutes les capsules n'acc\'el\`erent pas simultan\'ement), sous la tension $U_0 = \SI{72}{\volt}$ entre rails.

\begin{enumerate}[leftmargin=1.5em]
  \item En d\'eduire le courant moyen $I_{moy}$ circulant dans les rails.
  \item Le pont redresseur PD2 est constitu\'e de 4 diodes. Indiquer, en justifiant \`a partir du sens de conduction impos\'e par les diodes, si le bo\^itier d'alimentation est r\'eversible en courant du c\^ot\'e du r\'eseau alternatif.
  \item Lors du freinage d'une capsule, la machine synchrone qui l'entra\^ine fonctionne en g\'en\'eratrice et tend \`a renvoyer de l'\'energie \'electrique vers le bus continu des rails \`a $\SI{72}{\volt}$. Conclure sur le risque encouru si cette \'energie n'est pas g\'er\'ee autrement, compte tenu de la r\'eponse \`a la question pr\'ec\'edente, et indiquer bri\`evement une solution technique envisageable pour l'absorber.
\end{enumerate}

}{
\textbf{Corrig\'e.}

\begin{enumerate}[leftmargin=1.5em]
  \item $\langle i \rangle = I_{moy} = \dfrac{P_{moy}}{U_0} = \dfrac{7\,200}{72} = \SI{100}{\ampere}$.
  \item Dans un pont de diodes, chaque diode ne conduit que dans un seul sens (de l'anode vers la cathode) : le courant redress\'e $i_S(t)$ c\^ot\'e continu ne peut \^etre que positif ou nul, quel que soit le signe de la tension alternative d'entr\'ee. Le pont PD2 \`a diodes est donc \textbf{non r\'eversible en courant} : il ne peut transf\'erer de la puissance que du r\'eseau alternatif vers la charge continue, jamais dans l'autre sens.
  \item Puisque le pont PD2 ne peut pas laisser transiter un courant renvoy\'e vers le r\'eseau, l'\'energie de freinage r\'ecup\'er\'ee par la machine synchrone \textbf{ne peut pas \^etre \'evacu\'ee vers le r\'eseau EDF} \`a travers le bo\^itier d'alimentation : elle resterait pi\'eg\'ee sur le bus continu des rails, provoquant une surtension dangereuse pour les composants (filtre, hacheur, isolation TBTS) si rien ne l'absorbe. Une solution consiste \`a embarquer un dispositif de stockage local (par exemple des supercondensateurs dans la capsule) qui absorbe l'\'energie cin\'etique r\'ecup\'er\'ee lors du freinage, sans avoir besoin de la renvoyer \`a travers le bo\^itier d'alimentation non r\'eversible.
\end{enumerate}

\smallskip
\textbf{Conclusion.} Le courant de \SI{100}{\ampere} dans les rails et la non-r\'eversibilit\'e du pont \`a diodes fixent ensemble deux exigences de conception : un calibre de rail et de bo\^itier tenant \SI{100}{\ampere} en permanence, et un \'etage de dissipation ou de stockage embarqu\'e obligatoire, faute de quoi la tension du bus continu sortirait du domaine TBTS ($\leq \SI{120}{\volt}$ en continu) que le cahier des charges impose de respecter.
}
\end{tcolorbox}


% ============================================================
% EXERCICE 3 (niveau 2) -- Ondulation residuelle et nombre de pulsations
% ============================================================
\begin{tcolorbox}[exoBox,
  title={Exercice 3 -- \difficulte{2} \hfill\normalfont\footnotesize exercice in\'edit}
]
\CAPETSHOW{
On compare la qualit\'e de la tension continue obtenue en sortie de trois montages redresseurs \`a diodes aliment\'es par le m\^eme r\'eseau, en l'absence de tout filtrage : PD2 (monophas\'e, 2 pulsations), P3 (triphas\'e simple alternance, 3 pulsations) et PD6 (triphas\'e double alternance, 6 pulsations).

\medskip
\textbf{D\'efinition utilis\'ee dans tout l'exercice.} Le \emph{taux d'ondulation} $\tau$ est le rapport de la valeur \emph{efficace} $U_{ond,eff}$ de la composante alternative r\'esiduelle (l'ondulation) \`a la valeur moyenne de la tension redress\'ee :
\[
\tau = \frac{U_{ond,eff}}{\langle u_S \rangle}
\qquad\text{avec}\qquad
U_{ond,eff} = \sqrt{U_{S,eff}^{\,2} - \langle u_S \rangle^{2}}
\]
Les valeurs usuelles sont $\tau_{PD2} \approx 48{,}3\ \%$, $\tau_{P3} \approx 18{,}3\ \%$, $\tau_{PD6} \approx 4{,}2\ \%$. Cette d\'efinition ne doit pas \^etre confondue avec l'ondulation \emph{cr\^ete-\`a-cr\^ete} $\Delta U = U_{S,max} - U_{S,min}$, dont il sera question \`a la question 2.

\begin{center}
\begin{tikzpicture}[scale=0.85, >=Stealth]
\draw[-latex] (0,0) -- (9.3,0) node[right]{$\theta$};
\draw[-latex] (0,-0.4) -- (0,2.7) node[above]{$u_S(\theta)$};
\draw[dashed,colSlate] (0,2) -- (7.1,2) node[right,black,font=\small]{$U_{S,max}$};
\draw[dotted,colSlate] (0,1.27) -- (7.1,1.27) node[right,black,font=\small]{$\langle u_S\rangle$};
\foreach \k in {0,1,2}{%
  \draw[thick,colExo,domain=0:180,samples=60,smooth,variable=\x]
    plot ({\k*2.1+\x/180*2.1},{2*sin(\x)});}
\draw[latex-] (2.1,-0.06) -- (2.8,-0.85) node[right,font=\small]{$U_{S,min}$};
\node[font=\small] at (3.6,-1.45) {allure de $u_S(\theta)$ pour le montage PD2 non filtr\'e};
\end{tikzpicture}
\end{center}

\begin{enumerate}[leftmargin=1.5em]
  \item Justifier, sans calcul, pourquoi le taux d'ondulation diminue fortement lorsque le nombre de pulsations par p\'eriode augmente (PD2 $\to$ P3 $\to$ PD6).
  \item Un montage PD2 non filtr\'e alimente une charge r\'esistive dans les m\^emes conditions qu'\`a l'exercice 1 ($V_{eff} = \SI{230}{\volt}$, $f = \SI{50}{\hertz}$, sans filtrage).
  \begin{enumerate}[label=\alph*)]
    \item En reprenant les valeurs de $\langle u_S \rangle$ et de $U_{S,eff}$ \'etablies \`a l'exercice 1, calculer la valeur efficace $U_{ond,eff}$ de l'ondulation r\'esiduelle de deux mani\`eres : \`a partir de $\tau_{PD2}$, puis directement par la d\'efinition ci-dessus.
    \item Relever sur l'allure ci-dessus les valeurs extr\^emes $U_{S,max}$ et $U_{S,min}$, les calculer, et en d\'eduire l'ondulation cr\^ete-\`a-cr\^ete $\Delta U$.
    \item Comparer le rapport $\Delta U / \langle u_S \rangle$ au taux $\tau_{PD2}$ et commenter.
  \end{enumerate}
  \item Un cahier des charges impose un taux d'ondulation r\'esiduel $\tau$ inf\'erieur \`a $10\ \%$ sans filtrage suppl\'ementaire. Indiquer, parmi les trois montages, lequel (ou lesquels) satisfait ce crit\`ere.
  \item Le contenu harmonique du courant absorb\'e c\^ot\'e r\'eseau par un pont redresseur \`a diodes est constitu\'e des rangs $2k\pm1$ pour un PD2 et $6k\pm1$ pour un PD6 ($k$ entier). Expliquer en une phrase pourquoi un passage de PD2 \`a PD6 est favorable du point de vue de la compatibilit\'e \'electromagn\'etique c\^ot\'e r\'eseau.
\end{enumerate}

}{
\textbf{Corrig\'e.}

\begin{enumerate}[leftmargin=1.5em]
  \item Plus le nombre de pulsations par p\'eriode augmente, plus la tension redress\'ee est reconstitu\'ee \`a partir de portions de sinuso\"ide courtes et proches de son maximum (chaque portion couvre un angle $2\pi/q$ de plus en plus petit) : l'\'ecart entre la valeur instantan\'ee et la valeur moyenne se r\'eduit d'autant, donc l'ondulation r\'esiduelle diminue quand $q$ augmente.
  \item \begin{enumerate}[label=\alph*)]
    \item D'apr\`es l'exercice 1 : $\langle u_S \rangle = \dfrac{2\sqrt2}{\pi}V_{eff} \approx \SI{207}{\volt}$ et $U_{S,eff} = V_{eff} = \SI{230}{\volt}$.

    Par le taux d'ondulation : $U_{ond,eff} = \tau_{PD2}\times\langle u_S \rangle = 0{,}483\times 207 \approx \SI{100}{\volt}$.

    Directement par la d\'efinition :
    \[
    U_{ond,eff} = \sqrt{U_{S,eff}^{\,2} - \langle u_S \rangle^{2}} = \sqrt{230^2 - 207^2} = \sqrt{52\,900 - 42\,849} = \sqrt{10\,051} \approx \SI{100}{\volt}
    \]
    Les deux voies concordent, ce qui confirme bien $\tau_{PD2} = 100/207 \approx 0{,}483$.
    \item La tension redress\'ee est $u_S(\theta) = \sqrt2\,V_{eff}|\sin\theta|$ : son maximum est le sommet des arches et son minimum est atteint \`a chaque passage par z\'ero du r\'eseau, o\`u la tension redress\'ee \emph{s'annule} (aucun \'el\'ement de stockage ne maintient la tension en l'absence de filtrage) :
    \[
    U_{S,max} = \sqrt2\times 230 \approx \SI{325}{\volt}, \qquad U_{S,min} = \SI{0}{\volt}
    \]
    \[
    \Delta U = U_{S,max} - U_{S,min} \approx \SI{325}{\volt}
    \]
    \item $\dfrac{\Delta U}{\langle u_S \rangle} = \dfrac{325}{207} \approx 1{,}57$, soit $157\ \%$ : plus de trois fois la valeur de $\tau_{PD2} = 48{,}3\ \%$. Les deux grandeurs ne mesurent pas la m\^eme chose : $\tau$ rapporte la valeur \emph{efficace} de l'ondulation \`a la moyenne, alors que $\Delta U/\langle u_S \rangle$ en rapporte l'amplitude \emph{cr\^ete-\`a-cr\^ete}. M\'elanger les deux d\'efinitions (par exemple \'ecrire $\Delta U = \tau_{PD2}\times\langle u_S \rangle$) conduirait \`a une tension minimale de \SI{157}{\volt}, en contradiction directe avec l'allure trac\'ee, qui touche z\'ero deux fois par p\'eriode.
  \end{enumerate}
  \item Seul le montage PD6 ($\tau \approx 4{,}2\ \% < 10\ \%$) satisfait le crit\`ere sans filtrage suppl\'ementaire ; P3 ($18{,}3\ \%$) et PD2 ($48{,}3\ \%$) n\'ecessitent un filtre (LC par exemple) pour atteindre cette exigence.
  \item Le rang minimal des harmoniques inject\'es c\^ot\'e r\'eseau passe de $2k\pm1$ (soit le rang 1 -- fondamental -- puis 3, 5, \dots) \`a $6k\pm1$ (soit 5, 7, 11, 13, \dots pour un PD6) : les harmoniques de rang bas (le rang 3 notamment) disparaissent, ce qui r\'eduit le taux de distorsion harmonique global du courant r\'eseau et facilite le respect des normes de compatibilit\'e \'electromagn\'etique (par exemple la norme IEC/EN 61000-2-2, qui \guillemotleft\ fixe un THD de courant inf\'erieur \`a 25 \% pour assurer une compatibilit\'e \'electromagn\'etique conforme \guillemotright\ -- sujet CAPET SII 2024, \'epreuve \'ecrite disciplinaire, partie B).
\end{enumerate}

\smallskip
\textbf{Conclusion.} Avec une ondulation cr\^ete-\`a-cr\^ete \'egale \`a $157\ \%$ de sa valeur moyenne et une tension qui retombe \`a z\'ero deux fois par p\'eriode, un PD2 non filtr\'e ne peut pas alimenter directement une charge exigeant une tension quasi continue : le respect du crit\`ere $\tau < 10\ \%$ passe obligatoirement par un filtrage (cellule LC, condensateur de t\^ete) ou par le choix d'un montage \`a 6 pulsations.
}
\end{tcolorbox}


% ============================================================
% EXERCICE docTech (niveau 2) -- Urbanloop : lecture du DT1 (choix de section de
% rail) et chute de tension sur le bus continu 72 V issu du pont PD2
% ------------------------------------------------------------
% SOURCES (verifiees le 2026-09-02 dans E2_Appliquee_Commune_2024.md) :
%  - l.830  (DT1) : "Deformation de 0,13 mm pour un rail de dimensions 120x80x8 mm
%    et un espacement de 0,66 m" / "1,4 mm ... 120x80x8 mm ... 2 m" /
%    "0,08 mm ... 150x90x11 mm ... 0,66 m"  -> les 3 lignes du tableau docTech
%    sont le VERBATIM du DT1 (seule la mise en tableau est reconstituee, le DT1
%    d'origine etant une figure de courbes de simulation omise par la conversion).
%  - l.290  : "le cahier des charges impose une fleche maximale du rail inferieure
%    a 0,1 mm" (critere de la question 1)
%  - l.518  : "La resistivite de l'acier est rho = 0,13 Ohm.mm2.m-1"
%  - l.510  : "imposent en sortie une tension U0 de 72 V continue (TBTS)"
%  - l.498  : "Pmoy = 7,2 kW ... sous une tension U0 de 72 V"
%  - l.554  : "le cahier des charges impose une chute de tension maximale de 5% de
%    U0. Le courant moyen est pris egal a 100 A."
%  - l.217  : "la longueur par le rail est de 1,1 km entre la station 1 et la station 2"
%  - Q40 (l.540) : expression Req = 2.Rl.L.(x/L).(1-x/L) FOURNIE PAR LE SUJET
%  - RECONSTITUE : l'aire de la section (2519 mm2) n'est pas donnee numeriquement
%    par le sujet (figure 11 omise). Elle est ici RECALCULEE explicitement dans
%    l'enonce a partir des cotes 150x90x11 du DT1 en profil corniere en L, et la
%    demarche est demandee au candidat. Verification de coherence : cette section
%    conduit a Lmax = 1395 m > 1100 m, donc "deux boitiers suffisent", ce qui est
%    exactement la conclusion attendue par la question 43 du sujet.
% ============================================================
\begin{tcolorbox}[exoBox,
  title={Exercice 4 -- \difficulte{2} \hfill\normalfont\footnotesize 2024, \'Epreuve \'ecrite disciplinaire appliqu\'ee, questions 38 \`a 43 (DT1)}
]
\CAPETSHOW{
Les capsules du syst\`eme de transport Urbanloop sont aliment\'ees par des bo\^itiers r\'epartis le long de la voie. Chaque bo\^itier (pont PD2 + filtre + hacheur s\'erie) \guillemotleft\ impose en sortie une tension $U_0$ de \SI{72}{\volt} continue (TBTS : tr\`es basse tension de s\'ecurit\'e) entre les deux rails \guillemotright. Ce sont les \emph{rails eux-m\^emes} qui transportent ensuite la puissance jusqu'\`a la capsule : leur r\'esistance vient donc en s\'erie avec la sortie du redresseur et d\'egrade la tension effectivement vue par la capsule.

Le rail est un profil\'e en corni\`ere (section en L) en acier, de r\'esistivit\'e $\rho = \SI{0,13}{\ohm\milli\metre\squared\per\metre}$. Les r\'esultats de simulation m\'ecanique ayant servi \`a choisir ce profil\'e sont donn\'es ci-dessous.

\begin{tcolorbox}[docTech, title={Document technique DT1 -- r\'esultats de simulation (fl\`eche du rail)}]
\small
L'espacement indiqu\'e correspond \`a la distance entre deux traverses.

\smallskip
\begin{center}
\begin{tabular}{ccc}
\toprule
\textbf{Dimensions du rail} & \textbf{Espacement} & \textbf{D\'eformation (fl\`eche)} \\
\midrule
$120\times80\times8$\ \si{\milli\metre}  & \SI{0,66}{\metre} & \SI{0,13}{\milli\metre} \\
$120\times80\times8$\ \si{\milli\metre}  & \SI{2}{\metre}    & \SI{1,4}{\milli\metre}  \\
$150\times90\times11$\ \si{\milli\metre} & \SI{0,66}{\metre} & \SI{0,08}{\milli\metre} \\
\bottomrule
\end{tabular}
\end{center}

\smallskip
\emph{Exigence du cahier des charges m\'ecanique :} la fl\`eche maximale du rail doit rester inf\'erieure \`a \SI{0,1}{\milli\metre}.
\end{tcolorbox}

\smallskip
Un rail de corni\`ere de cotes $a \times b \times e$ (deux ailes de largeurs $a$ et $b$, d'\'epaisseur commune $e$) a pour aire de section $S = (a + b - e)\,e$.

\begin{enumerate}[leftmargin=1.5em]
  \item \textbf{Relever sur le document technique DT1} la seule configuration (dimensions de rail et espacement entre traverses) qui respecte l'exigence de fl\`eche du cahier des charges. Justifier en une phrase le rejet de chacune des deux autres.
  \item \`A partir des dimensions ainsi retenues, calculer l'aire $S$ de la section du rail, puis la r\'esistance \'electrique lin\'eique $R_\ell$ du rail, en \si{\ohm\per\metre}.
  \item La capsule se trouve \`a la distance $x$ du bo\^itier 1 et \`a la distance $L - x$ du bo\^itier 2, $L$ \'etant l'espacement entre deux bo\^itiers cons\'ecutifs. Les deux bo\^itiers alimentent la capsule en parall\`ele par l'interm\'ediaire des rails. Le sujet \'etablit que la r\'esistance \'equivalente vue par la capsule vaut
  \[ R_{eq} = 2\,R_\ell\,L\cdot\frac{x}{L}\left(1 - \frac{x}{L}\right) \]
  D\'eterminer la position $x$ la plus d\'efavorable au regard de la chute de tension, et en d\'eduire l'expression de la chute de tension maximale $\Delta u_{Rmax}$ en fonction de $R_\ell$, $L$ et du courant moyen $I_{moy}$.
  \item Le cahier des charges impose une chute de tension maximale de $5\ \%$ de $U_0$, le courant moyen \'etant pris \'egal \`a \SI{100}{\ampere}. D\'eterminer la distance maximale $L_{max}$ entre deux bo\^itiers d'alimentation. La longueur de voie \`a \'equiper entre les deux stations \'etant de \SI{1,1}{\kilo\metre}, v\'erifier que deux bo\^itiers d'alimentation suffisent pour ce parcours.
  \item Le pont PD2 du bo\^itier d\'elivre, avant filtrage et hachage, une tension moyenne $\langle u \rangle = \frac{2\sqrt2}{\pi}V_{eff}$ \`a partir du r\'eseau EDF monophas\'e $V_{eff} = \SI{230}{\volt}$. Calculer cette tension, puis le rapport cyclique du hacheur s\'erie qui en d\'eduit les \SI{72}{\volt} du bus. Indiquer si la chute de tension calcul\'ee \`a la question 4 remet en cause ce r\'eglage, et sous quelle condition.
\end{enumerate}

}{
\textbf{Corrig\'e.}

\begin{enumerate}[leftmargin=1.5em]
  \item \textbf{Lecture du DT1.} Seule la configuration \textbf{rail $150\times90\times11$\ \si{\milli\metre} avec un espacement de \SI{0,66}{\metre}} respecte l'exigence : sa fl\`eche vaut \SI{0,08}{\milli\metre} $<$ \SI{0,1}{\milli\metre}. Les deux autres sont rejet\'ees : le rail $120\times80\times8$ \`a \SI{0,66}{\metre} donne \SI{0,13}{\milli\metre} $>$ \SI{0,1}{\milli\metre}, et le m\^eme rail \`a \SI{2}{\metre} d'espacement donne \SI{1,4}{\milli\metre}, tr\`es largement hors crit\`ere (l'\'ecartement des traverses augmente fortement la fl\`eche).
  \item Aire de la section de la corni\`ere $150\times90\times11$ :
  \[ S = (a + b - e)\,e = (150 + 90 - 11)\times 11 = 229 \times 11 = \SI{2519}{\milli\metre\squared} \]
  R\'esistance lin\'eique (la r\'esistivit\'e \'etant donn\'ee en \si{\ohm\milli\metre\squared\per\metre}, elle se combine directement \`a une section en \si{\milli\metre\squared}) :
  \[ R_\ell = \frac{\rho}{S} = \frac{0{,}13}{2\,519} \approx \SI{5,16e-5}{\ohm\per\metre} = \SI{51,6}{\micro\ohm\per\metre} \]
  \item Posons $u = x/L \in [0\,;1]$. Alors $R_{eq} = 2R_\ell L \cdot u(1-u)$. La fonction $u \mapsto u(1-u)$ est maximale en $u = 1/2$, o\`u elle vaut $1/4$ : la position la plus d\'efavorable est donc la capsule \textbf{\`a mi-distance entre les deux bo\^itiers} ($x = L/2$), ce qui est physiquement attendu puisque c'est l\`a que les deux chemins de courant sont simultan\'ement les plus longs.
  \[ R_{eq,max} = 2R_\ell L \times \frac14 = \frac{R_\ell L}{2}
  \qquad\Longrightarrow\qquad
  \Delta u_{Rmax} = R_{eq,max}\,I_{moy} = \frac{R_\ell\,L\,I_{moy}}{2} \]
  \item Chute de tension admissible : $\Delta u_{adm} = 0{,}05 \times U_0 = 0{,}05 \times 72 = \SI{3,6}{\volt}$.
  \[ \frac{R_\ell\,L_{max}\,I_{moy}}{2} = \Delta u_{adm}
  \qquad\Longrightarrow\qquad
  L_{max} = \frac{2\,\Delta u_{adm}}{R_\ell\,I_{moy}} = \frac{2 \times 3{,}6}{5{,}16\times10^{-5} \times 100} \approx \SI{1395}{\metre} \]
  Deux bo\^itiers plac\'es aux deux extr\'emit\'es couvrent un espacement $L = \SI{1100}{\metre}$. Comme $1\,100 < L_{max} = 1\,395$, le crit\`ere est respect\'e : \textbf{deux bo\^itiers d'alimentation suffisent} pour les \SI{1,1}{\kilo\metre} de voie. La chute r\'eelle au point le plus d\'efavorable vaut alors
  \[ \Delta u_R = \frac{5{,}16\times10^{-5} \times 1\,100 \times 100}{2} \approx \SI{2,84}{\volt} \quad (3{,}9\ \% \text{ de } U_0) \]
  \item Tension moyenne en sortie du pont PD2 :
  \[ \langle u \rangle = \frac{2\sqrt2}{\pi}\times 230 \approx 0{,}9 \times 230 \approx \SI{207}{\volt} \]
  Le hacheur s\'erie (abaisseur) impose $\langle u_{sortie}\rangle = \alpha\,\langle u \rangle$, d'o\`u
  \[ \alpha = \frac{72}{207} \approx 0{,}35 \]
  Ce r\'eglage n'est pas remis en cause \`a condition que le hacheur soit \textbf{asservi} en tension : la chute de \SI{2,84}{\volt} appara\^it \emph{en aval} du bo\^itier, dans les rails, et non \`a ses bornes de sortie. Le bo\^itier maintient bien \SI{72}{\volt} \`a sa sortie ; c'est la capsule qui ne voit que \SI{69,2}{\volt} au point le plus d\'efavorable. Compenser cette chute en augmentant $\alpha$ (soit $\alpha \approx 0{,}36$ pour viser \SI{74,8}{\volt} en sortie de bo\^itier) rel\`everait la tension vue par une capsule proche d'un bo\^itier au-del\`a de \SI{74}{\volt} : la marge vis-\`a-vis de la limite TBTS (\SI{120}{\volt} en continu) reste confortable, mais le r\'eglage doit alors \^etre v\'erifi\'e \`a vide comme en charge.
\end{enumerate}

\smallskip
\textbf{Conclusion.} Le choix du profil\'e est ici doublement contraint : la section $150\times90\times11$\ \si{\milli\metre} est impos\'ee par la fl\`eche m\'ecanique (\SI{0,08}{\milli\metre} $<$ \SI{0,1}{\milli\metre}), et c'est cette m\^eme section qui, par sa faible r\'esistance lin\'eique (\SI{51,6}{\micro\ohm\per\metre}), autorise un espacement de \SI{1395}{\metre} entre bo\^itiers et rend donc suffisants les deux seuls bo\^itiers pr\'evus sur les \SI{1,1}{\kilo\metre} de voie. Un rail plus l\'eger aurait exig\'e un troisi\`eme bo\^itier d'alimentation : le dimensionnement m\'ecanique et le dimensionnement \'electrique de la distribution continue ne peuvent pas \^etre men\'es s\'epar\'ement.
}
\end{tcolorbox}

% ============================================================
% EXERCICE 5 (niveau 3) -- HYCITY, etude du P3 : commutateur plus positif,
% valeur moyenne par integration, courant efficace de diode, facteur de puissance
% ============================================================
\begin{tcolorbox}[exoBox,
  title={Exercice 5 -- \difficulte{3} \hfill\normalfont\footnotesize 2026, \'Epreuve 1 (admissibilit\'e), partie E, questions E2, E4, E6 et E8}
]
\CAPETSHOW{
Le bus \`a hydrog\`ene HYCITY\textregistered\ est \'equip\'e d'une batterie de secours recharg\'ee \`a partir du r\'eseau triphas\'e d'alimentation. La solution \'etudi\'ee ici associe le secondaire d'un transformateur triphas\'e et un pont redresseur \`a diodes de type P3. La charge, fortement inductive, est mod\'elis\'ee par une source de courant $I_S$ suppos\'e constant.

Les diodes sont \`a \textbf{cathodes communes} et suppos\'ees parfaites : le fonctionnement est celui d'un \emph{commutateur \guillemotleft\ plus positif \guillemotright}, c'est-\`a-dire que la diode voyant le plus haut potentiel \`a son anode est \`a l'\'etat passant, les autres \'etant bloqu\'ees.

\begin{center}
\begin{circuitikz}[european, scale=0.95]
% ---- point neutre (couplage etoile du secondaire) ----
\draw (-1.8,1.0) -- (-1.8,4.4);
\node[left] at (-1.9,2.7) {$N$};
% ---- trois sources (secondaire du transformateur) ----
\draw (-1.8,4.4) node[circ]{} to[sV, l=$v_1$] (0.9,4.4);
\draw (-1.8,2.7) node[circ]{} to[sV, l=$v_2$] (0.9,2.7);
\draw (-1.8,1.0) node[circ]{} to[sV, l=$v_3$] (0.9,1.0);
% ---- diodes : anode cote source, cathode en haut (cathodes communes) ----
\draw (0.9,4.4) to[D, l=$D_1$] (3.6,4.4);
\draw (0.9,2.7) to[D, l=$D_2$] (3.6,2.7);
\draw (0.9,1.0) to[D, l=$D_3$] (3.6,1.0);
% ---- barre des cathodes communes ----
\draw (3.6,1.0) -- (3.6,4.4);
\node[circ] at (3.6,2.7) {};
\draw (3.6,4.4) -- (5.6,4.4);
% ---- retour par le neutre ----
\draw (-1.8,1.0) -- (-1.8,-0.6) -- (5.6,-0.6);
% ---- charge : source de courant ----
\draw (5.6,4.4) to[I, l=$I_S$] (5.6,-0.6);
\draw[-latex] (7.1,4.4) -- (7.1,-0.6) node[midway,right]{$u_S$};
% ---- reperage de vD1 ----
\draw[-latex] (0.9,5.3) -- (3.6,5.3) node[midway,above,font=\small]{$v_{D1}$};
\draw[dotted] (0.9,4.5) -- (0.9,5.3);
\draw[dotted] (3.6,4.5) -- (3.6,5.3);
\end{circuitikz}
\end{center}

L'expression temporelle de la tension $v_1$ est $v_1(t) = \sqrt2\,V_{eff}\cos(\omega_1 t)$, avec $\omega_1 = 100\pi\ \si{\radian\per\second}$ la pulsation, $T_1$ la p\'eriode et $V_{eff}$ la valeur efficace des tensions simples au secondaire. Les tensions $v_2$ et $v_3$ sont d\'ephas\'ees respectivement de $-2\pi/3$ et $-4\pi/3$ par rapport \`a $v_1$.

\begin{enumerate}[leftmargin=1.5em]
  \item Dans le cas g\'en\'eral, \'etablir la relation liant $v_1$, $v_{D1}$ et $u_S$ (les tensions sont rep\'er\'ees sur le sch\'ema, le neutre $N$ servant de r\'ef\'erence). Simplifier cette relation lorsque la diode $D_1$ est \`a l'\'etat passant.
  \item Sur quel intervalle angulaire, en fonction de $\omega_1 t$, la diode $D_1$ est-elle passante ? En d\'eduire la p\'eriode $T_S$ de la tension $u_S$ en fonction de $T_1$.
  \item En utilisant la d\'efinition de la valeur moyenne d'une grandeur p\'eriodique
  \[ \langle x \rangle = \frac{1}{T}\int_{0}^{T} x(t)\,\mathrm dt \]
  montrer par le calcul int\'egral que la valeur moyenne de la tension redress\'ee s'exprime
  \[ \langle u_S \rangle = \frac{3\sqrt6}{2\pi}\,V_{eff} \]
  \item Repr\'esenter l'allure du courant $i_{D1}$ traversant la diode $D_1$ et le premier enroulement du secondaire, puis d\'eterminer, en fonction de $I_S$, la valeur efficace $I_{D1,eff}$ de ce courant.
  \item Le transformateur et les diodes \'etant parfaits, la puissance active est uniquement consomm\'ee par la charge et s'exprime $P_A = \langle u_S \rangle \cdot I_S$. Rappeler la d\'efinition du facteur de puissance $f_P$ au secondaire du transformateur, puis \'etablir sa valeur \`a partir des r\'esultats pr\'ec\'edents. Le cahier des charges impose $f_P > 0{,}9$ : conclure.
\end{enumerate}

}{
\textbf{Corrig\'e.}

\begin{enumerate}[leftmargin=1.5em]
  \item En parcourant la maille form\'ee par la source $v_1$, la diode $D_1$ et la charge (retour par le neutre $N$) :
  \[ v_1 = v_{D1} + u_S \qquad\text{soit}\qquad v_{D1} = v_1 - u_S \]
  Lorsque $D_1$ est passante, la diode est suppos\'ee parfaite donc $v_{D1} = 0$, d'o\`u
  \[ \boxed{u_S = v_1} \]
  Ce r\'esultat se g\'en\'eralise : \`a chaque instant, $u_S$ est \'egale \`a la plus grande des trois tensions simples, ce qui traduit exactement le fonctionnement en commutateur \guillemotleft\ plus positif \guillemotright.
  \item $D_1$ conduit tant que $v_1$ est sup\'erieure \`a $v_2$ et \`a $v_3$. Par sym\'etrie du syst\`eme triphas\'e \'equilibr\'e, les intersections de $v_1$ avec $v_2$ et $v_3$ encadrent le maximum de $v_1$ (atteint en $\omega_1 t = 0$ avec la convention en cosinus) \`a $\pm\pi/3$ :
  \[ D_1 \text{ passante pour } \omega_1 t \in \left[-\frac{\pi}{3},\ +\frac{\pi}{3}\right] \]
  soit un tiers de p\'eriode ($2\pi/3$ sur $2\pi$). Les trois diodes se relaient de fa\c{c}on identique et d\'ecal\'ee de $2\pi/3$ : le motif de $u_S$ se r\'ep\`ete donc trois fois par p\'eriode du r\'eseau,
  \[ \boxed{T_S = \frac{T_1}{3}} \]
  (montage \`a 3 pulsations par p\'eriode).
  \item Sur l'intervalle de conduction de $D_1$, $u_S(t) = v_1(t) = \sqrt2\,V_{eff}\cos(\omega_1 t)$. En int\'egrant sur une p\'eriode $T_S$ de $u_S$, c'est-\`a-dire sur $\omega_1 t \in [-\pi/3,\ +\pi/3]$ :
  \[
  \langle u_S \rangle = \frac{1}{2\pi/3}\int_{-\pi/3}^{+\pi/3} \sqrt2\,V_{eff}\cos\theta\;\mathrm d\theta
  = \frac{3\sqrt2\,V_{eff}}{2\pi}\Big[\sin\theta\Big]_{-\pi/3}^{+\pi/3}
  \]
  \[
  = \frac{3\sqrt2\,V_{eff}}{2\pi}\left(\frac{\sqrt3}{2} + \frac{\sqrt3}{2}\right)
  = \frac{3\sqrt2\,V_{eff}}{2\pi}\times\sqrt3
  = \frac{3\sqrt6}{2\pi}\,V_{eff}
  \]
  Num\'eriquement $\langle u_S \rangle \approx 1{,}17\,V_{eff}$ : la valeur moyenne redress\'ee est sup\'erieure \`a la valeur efficace d'une tension simple, car le montage ne pr\'el\`eve que le sommet des alternances.
  \item Le courant $I_S$ impos\'e par la charge inductive traverse int\'egralement l'enroulement dont la diode conduit. Le courant $i_{D1}$ est donc un cr\'eneau d'amplitude $I_S$ et de largeur $2\pi/3$ (soit un tiers de p\'eriode), nul le reste du temps :
  \begin{center}
  \begin{tikzpicture}[scale=0.85, >=Stealth]
  \draw[-latex] (-0.4,0) -- (8.6,0) node[right]{$\omega_1 t$};
  \draw[-latex] (0,-0.3) -- (0,2.2) node[above]{$i_{D1}$};
  % creneau 1 : de -pi/3 a pi/3 -> ici on centre sur 0 : trace partiel puis motif suivant
  \draw[thick,colExo] (0,1.6) -- (0.9,1.6) -- (0.9,0) -- (5.3,0) -- (5.3,1.6) -- (7.1,1.6) -- (7.1,0) -- (8.3,0);
  \draw[dashed,colSlate] (0,1.6) -- (5.3,1.6);
  \node[left,font=\small] at (0,1.6) {$I_S$};
  \draw[dotted] (0.9,0) -- (0.9,-0.55) node[below,font=\small]{$\pi/3$};
  \draw[dotted] (5.3,0) -- (5.3,-0.55) node[below,font=\small]{$5\pi/3$};
  \draw[dotted] (7.1,0) -- (7.1,-0.55) node[below,font=\small]{$7\pi/3$};
  \draw[latex-latex] (0.95,0.35) -- (5.25,0.35);
  \node[font=\footnotesize,fill=white,inner sep=1pt] at (3.1,0.35) {conduction de $D_2$ puis $D_3$};
  \end{tikzpicture}
  \end{center}
  Valeur efficace d'un cr\'eneau de rapport cyclique $1/3$ :
  \[
  I_{D1,eff} = \sqrt{\frac{1}{T_1}\int_{\text{conduction}} I_S^2\,\mathrm dt} = \sqrt{\frac{1}{3}\,I_S^2} = \frac{I_S}{\sqrt3} \approx 0{,}577\,I_S
  \]
  \item Le facteur de puissance au secondaire est le rapport de la puissance active \`a la puissance apparente :
  \[ f_P = \frac{P_A}{S} \qquad\text{avec}\qquad S = 3\,V_{eff}\,I_{D1,eff} \]
  (la puissance apparente d'un syst\`eme triphas\'e se calcule sur les valeurs \emph{efficaces} des tensions simples et des courants de ligne, ici les courants dans les enroulements secondaires).
  \[
  P_A = \langle u_S \rangle\,I_S = \frac{3\sqrt6}{2\pi}\,V_{eff} I_S \approx 1{,}17\,V_{eff} I_S
  \qquad
  S = 3\,V_{eff}\,\frac{I_S}{\sqrt3} = \sqrt3\,V_{eff} I_S \approx 1{,}73\,V_{eff} I_S
  \]
  \[
  f_P = \frac{3\sqrt6/(2\pi)}{\sqrt3} = \frac{3\sqrt2}{2\pi} \approx \mathbf{0{,}675}
  \]
  \textbf{Conclusion sur le cahier des charges.} $f_P = 0{,}675 < 0{,}9$ : le crit\`ere de facteur de puissance \textbf{n'est pas respect\'e} par le montage P3, et l'\'ecart au seuil est important (25\,\% en de\c{c}\`a). L'origine du d\'efaut est enti\`erement contenue dans les questions pr\'ec\'edentes : chaque enroulement ne conduit qu'un tiers de p\'eriode, si bien que la valeur efficace du courant qu'il supporte ($0{,}577\,I_S$) est disproportionn\'ee devant la puissance active effectivement transmise. La solution P3 doit donc \^etre \'ecart\'ee au profit d'une structure faisant conduire chaque enroulement plus longtemps (pont PD3).
\end{enumerate}

\smallskip
\textbf{Conclusion.} Le transformateur devrait \^etre dimensionn\'e pour une puissance apparente $S = 1{,}73\,V_{eff}I_S$ alors que la charge n'en consomme que $1{,}17\,V_{eff}I_S$ : pr\`es de 48\,\% de puissance apparente install\'ee en exc\`es, factur\'ee au constructeur en cuivre et en fer sans contrepartie utile. C'est bien un crit\`ere de dimensionnement, et non un simple indicateur de qualit\'e r\'eseau, qui condamne ici le P3 face \`a l'exigence $f_P > 0{,}9$ du cahier des charges.
}
\end{tcolorbox}

% ============================================================
% EXERCICE 6 (niveau 3) -- PD2 tout thyristor, commande de l'angle d'amorcage
% ============================================================
\begin{tcolorbox}[exoBox,
  title={Exercice 6 -- \difficulte{3} \hfill\normalfont\footnotesize exercice in\'edit}
]
\CAPETSHOW{
Un pont de Gra\"etz command\'e (PD2 tout thyristor) alimente une charge fortement inductive, mod\'elis\'ee par une source de courant constant $I_S$. La tension d'entr\'ee est $v(t) = \sqrt2\,V_{eff}\sin(\omega t)$. Les quatre thyristors sont amorc\'es avec un retard angulaire $\alpha$ (angle d'amor\c cage) par rapport au passage par z\'ero de $v(t)$ : la paire $T_1$-$T_2$ est amorc\'ee \`a $\theta = \alpha$ et conduit jusqu'\`a $\theta = \pi + \alpha$, o\`u la paire $T_3$-$T_4$ prend le relais.

\begin{center}
\begin{tikzpicture}[scale=0.85, >=Stealth]
\draw[-latex] (0,0) -- (7.9,0) node[right]{$\theta$};
\draw[-latex] (0,-1.95) -- (0,2.4) node[above]{$u_S(\theta)$};
\draw[dashed,colSlate] (0,1.8) -- (7.4,1.8);
\node[left,font=\small] at (0,1.8) {$\sqrt2 V_{eff}$};
% arc 1 : theta de alpha a pi+alpha (alpha = 30 deg)
\draw[thick,colExo,domain=30:210,samples=90,smooth,variable=\x]
  plot ({\x*0.01778},{1.8*sin(\x)});
% arc 2 : motif identique decale de pi (periode de repetition = pi)
\draw[thick,colExo,domain=30:210,samples=90,smooth,variable=\x]
  plot ({(\x+180)*0.01778},{1.8*sin(\x)});
% saut de commutation en pi+alpha
\draw[thick,colExo] (3.734,-0.9) -- (3.734,0.9);
\draw[latex-] (3.79,0) -- (4.55,-1.05) node[right,font=\footnotesize]{saut de commutation};
\draw[dotted] (0.533,0) -- (0.533,-1.55) node[below,font=\small]{$\alpha$};
\draw[dotted] (3.734,-0.9) -- (3.734,-1.55) node[below,font=\small]{$\pi+\alpha$};
\draw[dotted] (6.934,-0.9) -- (6.934,-1.55) node[below,font=\small]{$2\pi+\alpha$};
\end{tikzpicture}
\end{center}

\begin{enumerate}[leftmargin=1.5em]
  \item Justifier pourquoi, la charge \'etant fortement inductive (courant $I_S$ suppos\'e constant, jamais nul), une seule paire de thyristors conduit \`a la fois et pourquoi $u_S(\theta) = \sqrt2\,V_{eff}\sin\theta$ sur l'intervalle de conduction $[\alpha,\ \pi+\alpha]$.
  \item Expliquer l'origine du saut de tension observ\'e sur le chronogramme en $\theta = \pi+\alpha$ et donner son amplitude en fonction de $V_{eff}$ et $\alpha$.
  \item Montrer, par le calcul de la valeur moyenne de $u_S(\theta)$ sur une p\'eriode de conduction, que
  \[ \langle u_S \rangle = \underbrace{\frac{2\sqrt2}{\pi}V_{eff}}_{\langle u_S \rangle_0}\cos\alpha \]
  \item Application num\'erique : pour $V_{eff} = \SI{230}{\volt}$, calculer $\langle u_S \rangle_0$ (angle d'amor\c cage nul). Calculer $\langle u_S \rangle$ pour $\alpha = \ang{30}$. Pour quel angle $\alpha$ a-t-on $\langle u_S \rangle = 0$ ? Que peut-on dire du sens de transfert de la puissance pour $\alpha > \ang{90}$ ?
\end{enumerate}

}{
\textbf{Corrig\'e.}

\begin{enumerate}[leftmargin=1.5em]
  \item Le courant $I_S$ ne peut jamais s'annuler brutalement dans une charge fortement inductive : \`a tout instant, un et un seul couple de thyristors doit \^etre passant pour assurer la continuit\'e de $I_S$. Tant que $T_1$-$T_2$ sont amorc\'es et passants (de $\alpha$ \`a $\pi+\alpha$), la tension aux bornes de la charge est directement impos\'ee par la source : $u_S(\theta) = v(\theta) = \sqrt2\,V_{eff}\sin\theta$ (thyristors suppos\'es parfaits, chute de tension nulle \`a l'\'etat passant).
  \item En $\theta = (\pi+\alpha)^-$, $T_1$-$T_2$ conduisent encore et $u_S = \sqrt2\,V_{eff}\sin(\pi+\alpha) = -\sqrt2\,V_{eff}\sin\alpha$ (valeur \textbf{n\'egative}). \`A $\theta = \pi+\alpha$, l'amor\c cage de $T_3$-$T_4$ commute la charge sur l'autre diagonale du pont : la source est alors vue \`a l'envers, donc $u_S = -v(\theta) = -\sqrt2\,V_{eff}\sin(\pi+\alpha) = +\sqrt2\,V_{eff}\sin\alpha$. La tension saute donc instantan\'ement de $-\sqrt2\,V_{eff}\sin\alpha$ \`a $+\sqrt2\,V_{eff}\sin\alpha$, soit une amplitude de saut
  \[ \Delta u_S = 2\sqrt2\,V_{eff}\sin\alpha \]
  Le motif se r\'ep\`ete ensuite \`a l'identique : la p\'eriode de $u_S$ vaut $\pi$ (2 pulsations par p\'eriode du r\'eseau), et non $2\pi$.
  \item Valeur moyenne sur la p\'eriode de conduction $[\alpha,\ \pi+\alpha]$ (de largeur $\pi$, qui constitue la p\'eriode de r\'ep\'etition du motif $u_S$) :
  \[
  \langle u_S \rangle = \frac1\pi\int_\alpha^{\pi+\alpha} \sqrt2\,V_{eff}\sin\theta\,\mathrm d\theta
  = \frac{\sqrt2\,V_{eff}}{\pi}\Big[-\cos\theta\Big]_\alpha^{\pi+\alpha}
  = \frac{\sqrt2\,V_{eff}}{\pi}\big(\cos\alpha - \cos(\pi+\alpha)\big)
  \]
  Or $\cos(\pi+\alpha) = -\cos\alpha$, donc :
  \[
  \langle u_S \rangle = \frac{\sqrt2\,V_{eff}}{\pi}\times 2\cos\alpha = \frac{2\sqrt2}{\pi}V_{eff}\cos\alpha
  \]
  On retrouve bien $\langle u_S \rangle = \langle u_S \rangle_0 \cos\alpha$ avec $\langle u_S \rangle_0 = \dfrac{2\sqrt2}{\pi}V_{eff}$, la valeur moyenne du pont PD2 \`a diodes (montage non command\'e, $\alpha = 0$).
  \item $\langle u_S \rangle_0 = \dfrac{2\sqrt2}{\pi}\times 230 \approx 0{,}9\times 230 \approx \SI{207}{\volt}$.

  Pour $\alpha = \ang{30}$ : $\langle u_S \rangle = 207\times\cos(\ang{30}) = 207\times 0{,}866 \approx \SI{179}{\volt}$.

  $\langle u_S \rangle = 0$ pour $\cos\alpha = 0$, soit $\alpha = \ang{90}$.

  Pour $\alpha > \ang{90}$, $\cos\alpha < 0$ donc $\langle u_S \rangle < 0$ : la tension moyenne aux bornes de la charge s'inverse alors que le courant $I_S$ garde le m\^eme sens (impos\'e par la charge inductive). La puissance moyenne $\langle u_S \rangle \times I_S$ devient n\'egative : le montage transf\`ere alors de la puissance de la charge continue vers le r\'eseau alternatif -- on parle de fonctionnement en \textbf{onduleur assist\'e} (ou onduleur non autonome), possible uniquement gr\^ace \`a la commandabilit\'e des thyristors (impossible avec un pont \`a diodes, non commandable).
\end{enumerate}

\smallskip
\textbf{Conclusion.} Le r\'eglage de $\alpha$ entre \ang{0} et \ang{180} donne acc\`es \`a une tension moyenne continu\^ument variable de $+\SI{207}{\volt}$ \`a $-\SI{207}{\volt}$ avec un courant de charge de signe constant : le pont tout thyristor couvre donc les quadrants 1 et 4 du plan $(\langle u_S \rangle,\ I_S)$, ce qui permet le freinage par r\'ecup\'eration -- exigence qu'un pont \`a diodes, limit\'e au seul quadrant 1, ne peut pas satisfaire (cf. exercice 2).
}
\end{tcolorbox}

\CAPETENDFILE
